\documentclass[11pt,a4paper]{article}
\usepackage[utf8]{inputenc}
\usepackage[T1]{fontenc}
\IfFileExists{french.ldf}{\usepackage[french]{babel}}{\usepackage[english]{babel}}
\usepackage{amsmath,amssymb,amsthm}
\usepackage{geometry}\geometry{margin=2.3cm}
\usepackage{listings}
\usepackage{xcolor}
\usepackage{enumitem}
\usepackage{fancyhdr}
\usepackage{lmodern}
\usepackage{titlesec}
\usepackage{hyperref}
\hypersetup{colorlinks=true, urlcolor=[RGB]{2,101,115}, linkcolor=black}
\definecolor{codebg}{RGB}{247,249,251}
\definecolor{kw}{RGB}{175,45,45}
\definecolor{cmt}{RGB}{110,120,130}
\definecolor{str}{RGB}{20,110,60}
\lstset{language=Python, backgroundcolor=\color{codebg}, basicstyle=\ttfamily\small,
  keywordstyle=\color{kw}\bfseries, commentstyle=\color{cmt}\itshape, stringstyle=\color{str},
  numbers=left, numberstyle=\tiny\color{cmt}, numbersep=8pt, frame=single, rulecolor=\color{gray!40},
  showstringspaces=false, breaklines=true, columns=fullflexible, extendedchars=true, inputencoding=utf8,
  literate={é}{{\'e}}1 {è}{{\`e}}1 {à}{{\`a}}1 {ê}{{\^e}}1 {î}{{\^i}}1 {ç}{{\c c}}1 {ô}{{\^o}}1 {û}{{\^u}}1 {ù}{{\`u}}1 {É}{{\'E}}1}
\newcounter{qcnt}
\newcommand{\Q}{\medskip\par\stepcounter{qcnt}\noindent\textbf{Question~\theqcnt.}~}
\newcommand{\code}[1]{\lstinline!#1!}
\pagestyle{fancy}\fancyhf{}
\lhead{\small Préparation au CNC 2026 --- Informatique MP}
\rhead{\small Algorithmique de la blockchain Ethereum}
\cfoot{\thepage}
\renewcommand{\headrulewidth}{0.4pt}
\titleformat*{\section}{\large\bfseries\scshape}

\begin{document}

\begin{center}
{\large\bfseries Préparation au Concours National Commun 2026 --- Informatique MP}\\[8pt]
\rule{0.9\linewidth}{0.5pt}\\[6pt]
{\Large\bfseries Épreuve d'Informatique}\\[4pt]
{\bfseries Filière MP --- Durée : 4 heures}\\[6pt]
{\itshape Algorithmique de la blockchain Ethereum : gaz, encodage, arbres de Merkle, aléa et programmation dynamique}\\[6pt]
\rule{0.9\linewidth}{0.5pt}\\[8pt]
{\small Auteur~: \href{https://orcid.org/0000-0002-6108-7176}{\textbf{Pr Agrégé El Hadiq Zouhair}}}\\[3pt]
{\small \href{https://cpgeacademy.org}{\textbf{cpgeacademy.org}} \quad---\quad Tous droits réservés \textcopyright\ cpgeacademy.org}
\end{center}
\medskip
\begin{center}
\fcolorbox{gray}{gray!5}{\begin{minipage}{0.93\linewidth}\small
\textbf{Avis important.} Ce document est une \textbf{initiative personnelle} du professeur,
à but strictement pédagogique. Il ne s'agit \textbf{pas} d'un sujet officiel du Concours
National Commun~; il s'en inspire uniquement par la forme et l'esprit.
\smallskip
\textbf{Pour aller plus loin.} Les élèves peuvent traiter une \textbf{nouvelle version} de
ce problème dotée d'un \textbf{autocorrecteur des réponses}, et suivre un \textbf{cours
complet}, sur \href{https://cpgeacademy.org}{\textbf{cpgeacademy.org}}.
\end{minipage}}
\end{center}

\bigskip
\section*{Présentation et notations}

Ethereum est un ordinateur mondial décentralisé. Ce problème étudie, de manière \emph{largement
indépendante} partie par partie, quelques structures et algorithmes qui le sous-tendent~: la
mesure des ressources (le \emph{gaz}), l'encodage hexadécimal des adresses, les \emph{arbres de
Merkle}, la sélection aléatoire d'un validateur, et enfin un problème d'optimisation résolu par
\emph{programmation dynamique}.

\noindent\textbf{Conventions.}
\begin{itemize}[nosep]
  \item Le langage est Python 3. Les entiers Python sont de taille arbitraire.
  \item Pour une liste \code{L}, \code{len(L)} est sa longueur et \code{L[i]} son terme d'indice $i$ (à partir de $0$).
  \item $\lfloor x \rfloor$ désigne la partie entière de $x$. $\log_b$ est le logarithme en base $b$.
  \item On rappelle que \code{n // b} et \code{n \% b} sont le quotient et le reste de la division euclidienne.
\end{itemize}

\noindent Les questions sont numérotées \textbf{en continu} sur tout le sujet.

\bigskip
\hrule
\bigskip
{\Large\bfseries É\,N\,O\,N\,C\,É}
\medskip

\section*{Partie I --- Unités et compteur de gaz}

L'unité de compte d'Ethereum est l'\emph{ether}~; sa plus petite subdivision est le \emph{wei},
avec $1$ ether $= 10^{18}$ wei. Chaque instruction exécutée consomme une quantité fixe de
\emph{gaz}~; une transaction fixe un plafond de gaz au-delà duquel l'exécution est interrompue.

\Q Écrire une fonction \code{ether_vers_wei(e)} qui, pour un entier $e \ge 0$, renvoie le nombre
de wei correspondant à $e$ ethers.

\Q Écrire \code{gaz_restant(plafond, couts)} où \code{couts} est la liste des coûts (entiers
positifs) des instructions exécutées dans l'ordre. La fonction retranche les coûts un à un~; si,
avant une instruction, le gaz disponible est \emph{strictement} inférieur à son coût, l'exécution
s'arrête et la fonction renvoie la chaîne \code{"panne de gaz"}. Sinon elle renvoie le gaz restant.

\Q Un contrat exécute une boucle dont l'itération numéro $i$ (pour $i \ge 0$) coûte $g/2^{i}$
unités de gaz, où $g$ est un réel strictement positif. Montrer que le coût total
$\displaystyle C = \sum_{i=0}^{+\infty} \frac{g}{2^{i}}$ est fini et vaut exactement $2g$. En
déduire que, quel que soit le nombre d'itérations effectivement exécutées, le coût cumulé est
majoré par $2g$.

\section*{Partie II --- Encodage hexadécimal des adresses}

Une adresse Ethereum est un entier affiché en base $16$ (hexadécimal), avec l'alphabet
\code{"0123456789abcdef"}. On étudie l'encodage d'un entier en base $16$.

\Q Écrire \code{entier_vers_hex(n)} qui, pour un entier $n \ge 0$, renvoie la chaîne de ses
chiffres hexadécimaux (sans le préfixe \code{0x}), avec la convention \code{entier_vers_hex(0) == "0"}.

\Q Écrire \code{hex_vers_entier(s)} qui, pour une chaîne \code{s} de chiffres hexadécimaux
minuscules, renvoie l'entier qu'elle représente.

\Q Pour $n \ge 1$, montrer que le nombre de tours de la boucle de \code{entier_vers_hex(n)}
(donc le nombre de chiffres produits) est exactement $\lfloor \log_{16} n \rfloor + 1$. En déduire
la complexité temporelle de \code{entier_vers_hex} en fonction du nombre de chiffres.

\Q Justifier la \textbf{terminaison} de la boucle de \code{entier_vers_hex} à l'aide d'un
\emph{variant de boucle}, puis prouver la \textbf{correction}, c'est-à-dire que pour tout
$n \ge 0$, \code{hex_vers_entier(entier_vers_hex(n))} est égal à $n$.

\section*{Partie III --- Arbres de Merkle}

Un \emph{arbre de Merkle} résume un ensemble de données par un seul haché, la \emph{racine}. On
se donne une fonction de combinaison
$$\mathtt{parent(a, b)} = (31\,a + b) \bmod 1000003,$$
et les feuilles sont des entiers. La racine se calcule par niveaux~: à chaque niveau, on remplace
la liste courante par la liste des \code{parent} de couples consécutifs~; si la liste a une
longueur impaire, on \textbf{duplique son dernier élément} avant de l'appairer. On s'arrête
lorsqu'il ne reste qu'un élément.

\Q Écrire \code{racine_merkle(feuilles)} qui prend une liste non vide de feuilles et renvoie la
racine. Vérifier à la main que \code{racine_merkle([1, 2, 3, 4])} vaut $1120$.

\Q Soit $n$ le nombre de feuilles. Montrer que le nombre total d'appels à \code{parent} effectués
par \code{racine_merkle} est un $O(n)$. \emph{Indication}~: majorer la somme des tailles des
niveaux successifs.

\Q On veut prouver l'\emph{inclusion} d'une feuille sans révéler tout l'arbre. Écrire
\code{preuve_inclusion(feuilles, i)} qui renvoie la liste des couples $(\text{frère}, \text{côté})$
rencontrés en remontant de la feuille d'indice $i$ jusqu'à la racine, où le côté vaut
\code{'L'} ou \code{'R'} et indique si le frère est à gauche ou à droite. Écrire ensuite
\code{verifier(feuille, preuve, racine)} qui recalcule la racine à partir de la feuille et de la
preuve, et renvoie un booléen.

\Q Montrer par récurrence sur la hauteur de l'arbre que, si la fonction \code{parent} est
\emph{résistante aux collisions} (on ne sait pas produire $ (a,b) \ne (a',b')$ avec
$\mathtt{parent}(a,b) = \mathtt{parent}(a',b')$), alors une preuve d'inclusion acceptée par
\code{verifier} garantit que la feuille figure bien parmi les feuilles ayant produit la racine.
Quelle est la taille d'une preuve d'inclusion en fonction de $n$~?

\section*{Partie IV --- Sélection aléatoire d'un validateur}

À chaque créneau (\emph{slot}), un validateur donné est choisi comme proposeur avec probabilité
$p \in\, ]0,1[$, indépendamment des autres créneaux. On note $N$ le numéro (à partir de $1$) du
premier créneau où ce validateur est choisi.

\Q Donner, pour $k \ge 1$, la probabilité $\mathbb{P}(N = k)$. Reconnaître la loi de $N$ et en
déduire l'espérance $\mathbb{E}[N]$.

\Q Montrer que $\mathbb{P}(N > k) = (1-p)^{k}$, puis que $\mathbb{P}(N < +\infty) = 1$
(\emph{terminaison presque sûre}~: le validateur est choisi en temps fini avec probabilité $1$).

\Q Application numérique~: pour $p = 1/4$, donner $\mathbb{E}[N]$ et la probabilité
$\mathbb{P}(N \le 4)$ sous forme de fraction irréductible.

\section*{Partie V --- Partition équilibrée par programmation dynamique}

On dispose d'une liste $S$ d'entiers strictement positifs. Le \emph{problème de partition}
demande s'il existe un sous-ensemble de $S$ dont la somme vaut exactement la moitié de la somme
totale (les deux parts ont alors la même somme). C'est un problème classique lié aux preuves à
divulgation nulle.

\Q Montrer que si la somme totale $T = \sum S$ est impaire, la réponse est nécessairement
\emph{non}. On suppose désormais $T$ paire et on pose $C = T/2$.

\Q Le \emph{problème de décision} associé (« existe-t-il un sous-ensemble de somme $C$~? ») est
\textbf{NP}. Justifier en décrivant un \emph{certificat} et l'algorithme de vérification, ainsi
que sa complexité.

\Q On définit, pour $0 \le i \le \mathtt{len(S)}$ et $0 \le c \le C$, le booléen $A[i][c]$ = « il
existe un sous-ensemble des $i$ premiers éléments de $S$ de somme $c$ ». Donner la relation de
récurrence vérifiée par $A[i][c]$ et son initialisation.

\Q En déduire une fonction \code{peut_partitionner(S)} utilisant un tableau à une dimension
\code{dp} de taille $C+1$. \emph{Attention au sens de parcours} pour ne pas réutiliser deux fois
le même élément. Donner sa complexité en temps et en espace, et expliquer le terme
\emph{pseudo-polynomial}.

\Q Prouver la correction de \code{peut_partitionner} en énonçant un \emph{invariant} portant sur
\code{dp} après le traitement des $i$ premiers éléments, puis en le démontrant par récurrence sur
$i$. Justifier enfin la terminaison.

\newpage
\setcounter{qcnt}{0}
\bigskip
\hrule
\bigskip
{\Large\bfseries C\,O\,R\,R\,I\,G\,É}
\medskip

\section*{Partie I --- Unités et compteur de gaz}

\Q On multiplie par $10^{18}$.
\begin{lstlisting}
def ether_vers_wei(e):
    return e * 10**18
\end{lstlisting}
Complexité $O(1)$ opérations sur grands entiers.

\Q On parcourt les coûts en maintenant le gaz disponible.
\begin{lstlisting}
def gaz_restant(plafond, couts):
    restant = plafond
    for c in couts:
        if restant < c:
            return "panne de gaz"
        restant -= c
    return restant
\end{lstlisting}
La boucle est finie (\code{len(couts)} tours), d'où la terminaison~; complexité $O(m)$ avec
$m = \mathtt{len(couts)}$.

\Q C'est une série géométrique de raison $1/2 \in\, ]0,1[$, donc convergente~:
$$ C = \sum_{i=0}^{+\infty} g\Bigl(\tfrac12\Bigr)^{i} = \frac{g}{1 - \tfrac12} = 2g. $$
Une exécution qui s'arrête après $m$ itérations consomme
$\sum_{i=0}^{m-1} g/2^{i} = g\,\dfrac{1-(1/2)^{m}}{1-1/2} = 2g\,(1 - 2^{-m}) < 2g.$
Le coût cumulé est donc \emph{toujours} strictement majoré par $2g$, indépendamment de $m$.
$\square$

\section*{Partie II --- Encodage hexadécimal des adresses}

\Q
\begin{lstlisting}
def entier_vers_hex(n):
    if n == 0:
        return "0"
    chiffres = "0123456789abcdef"
    s = ""
    while n > 0:
        s = chiffres[n % 16] + s
        n //= 16
    return s
\end{lstlisting}

\Q
\begin{lstlisting}
def hex_vers_entier(s):
    chiffres = "0123456789abcdef"
    n = 0
    for c in s:
        n = n * 16 + chiffres.index(c)
    return n
\end{lstlisting}

\Q Soit $n \ge 1$ et $d$ le nombre de chiffres produits. À chaque tour, $n$ est remplacé par
$\lfloor n/16 \rfloor$~; la boucle s'arrête dès que $n$ atteint $0$. Le nombre de tours est donc
le plus petit $d$ tel que $\lfloor n / 16^{d} \rfloor = 0$, c'est-à-dire $16^{d} > n$, soit
$d > \log_{16} n$. Le plus petit tel entier est $d = \lfloor \log_{16} n \rfloor + 1$. Chaque tour
coûte $O(1)$ (hors coût des grands entiers), donc la complexité est $\Theta(d) = \Theta(\log n)$
tours.

\Q \textbf{Terminaison.} La quantité entière $n \ge 0$ décroît strictement à chaque tour
(on passe de $n$ à $\lfloor n/16 \rfloor < n$ tant que $n \ge 1$) et reste positive~: c'est un
\emph{variant de boucle}, donc la boucle termine.

\noindent\textbf{Correction.} Notons $c_{d-1}\dots c_1 c_0$ les chiffres produits (poids faible
$c_0$). Par construction, à chaque tour on extrait $c_j = n_j \bmod 16$ puis $n_{j+1} =
\lfloor n_j / 16\rfloor$, avec $n_0 = n$. Par la division euclidienne, $n_j = 16\,n_{j+1} + c_j$
avec $0 \le c_j < 16$. Une récurrence immédiate donne $n = \sum_{j=0}^{d-1} c_j\, 16^{j}$. Or
\code{hex_vers_entier} évalue par schéma de Horner exactement $\sum_{j=0}^{d-1} c_j\,16^{j}$. Donc
\code{hex_vers_entier(entier_vers_hex(n))} est égal à $n$. Le cas $n=0$ est traité à part et renvoie
\code{"0"}, décodé en $0$. $\square$

\medskip
\noindent\emph{Vérification numérique}~: $255 \mapsto \mathtt{"ff"}$, $16 \mapsto \mathtt{"10"}$,
$7271972807 \mapsto \mathtt{"1b1717fc7"}$ ($9$ chiffres, et $\lfloor\log_{16}
7271972807\rfloor+1 = 9$).

\section*{Partie III --- Arbres de Merkle}

\Q
\begin{lstlisting}
def parent(a, b):
    return (31 * a + b) % 1000003

def racine_merkle(feuilles):
    niveau = list(feuilles)
    while len(niveau) > 1:
        if len(niveau) % 2 == 1:
            niveau.append(niveau[-1])
        niveau = [parent(niveau[i], niveau[i+1])
                  for i in range(0, len(niveau), 2)]
    return niveau[0]
\end{lstlisting}
Vérification~: niveau $0 = [1,2,3,4]$~; $\mathtt{parent}(1,2)=33$, $\mathtt{parent}(3,4)=97$~; niveau
$1 = [33,97]$~; $\mathtt{parent}(33,97)=33\cdot31+97=1120$. Donc la racine vaut $\mathbf{1120}$.
(De même $\mathtt{racine\_merkle}([1,2,3])=1119$ après duplication du $3$.)

\Q À chaque niveau, si la liste a une taille $t$, on effectue $\lceil t/2\rceil$ appels à
\code{parent} et la taille suivante est $\lceil t/2\rceil$. Partant de $t_0 = n$, on a
$t_{k+1} = \lceil t_k/2\rceil \le t_k/2 + 1$. Le nombre total d'appels est
$\sum_{k\ge 0} \lceil t_k/2\rceil = \sum_{k\ge 1} t_k \le \sum_{k\ge 1}\bigl(n/2^{k} + 2\bigr)$.
La partie géométrique se majore par $n$, et il y a $O(\log n)$ niveaux, donc le total est
$\le n + O(\log n) = O(n)$. $\square$

\Q
\begin{lstlisting}
def preuve_inclusion(feuilles, i):
    niveau = list(feuilles)
    preuve = []
    idx = i
    while len(niveau) > 1:
        if len(niveau) % 2 == 1:
            niveau.append(niveau[-1])
        if idx % 2 == 0:
            preuve.append((niveau[idx+1], 'R'))
        else:
            preuve.append((niveau[idx-1], 'L'))
        niveau = [parent(niveau[j], niveau[j+1])
                  for j in range(0, len(niveau), 2)]
        idx //= 2
    return preuve

def verifier(feuille, preuve, racine):
    acc = feuille
    for frere, cote in preuve:
        if cote == 'R':
            acc = parent(acc, frere)
        else:
            acc = parent(frere, acc)
    return acc == racine
\end{lstlisting}
Vérification~: $\mathtt{preuve\_inclusion}([1,2,3,4],0) = [(2,\mathtt{'R'}),(97,\mathtt{'R'})]$~; puis
$\mathtt{parent}(1,2)=33$, $\mathtt{parent}(33,97)=1120=$ racine.

\Q \textbf{Récurrence sur la hauteur $h$.} Pour $h=0$ (une seule feuille), la racine \emph{est}
la feuille~; une preuve vide n'est acceptée que si $\mathtt{feuille} = \mathtt{racine}$. Supposons la
propriété vraie à la hauteur $h$ et considérons un arbre de hauteur $h+1$. La vérification
recalcule d'abord $r' = \mathtt{parent}(\mathtt{feuille}, \text{frère})$ (ou l'inverse), puis vérifie
que $r'$, muni du reste de la preuve, mène à la racine~: c'est une preuve d'inclusion de $r'$
dans un arbre de hauteur $h$, valide par hypothèse de récurrence, donc $r'$ est bien un n\oe ud
interne du niveau $h$. Comme \code{parent} est résistante aux collisions, l'unique antécédent
connu de $r'$ est le couple $(\mathtt{feuille}, \text{frère})$ réellement présent~; la feuille figure
donc parmi celles ayant produit la racine. $\square$

\noindent La preuve contient un couple par niveau, soit une taille $O(\log n)$~: c'est
l'intérêt majeur des arbres de Merkle (certificat logarithmique).

\section*{Partie IV --- Sélection aléatoire d'un validateur}

\Q Le premier succès arrive au $k$-ième créneau si les $k-1$ premiers sont des échecs
(probabilité $(1-p)^{k-1}$) suivis d'un succès (probabilité $p$)~:
$$\mathbb{P}(N = k) = (1-p)^{k-1}\,p, \qquad k \ge 1.$$
C'est la \emph{loi géométrique} de paramètre $p$, d'espérance
$\mathbb{E}[N] = \sum_{k\ge1} k (1-p)^{k-1} p = \dfrac{1}{p}.$

\Q $\mathbb{P}(N > k)$ est la probabilité de $k$ échecs consécutifs, soit $(1-p)^{k}$. Comme
$0 < 1-p < 1$, $(1-p)^{k} \to 0$ quand $k \to +\infty$, donc
$\mathbb{P}(N < +\infty) = 1 - \lim_k \mathbb{P}(N>k) = 1$~: terminaison presque sûre. $\square$

\Q Pour $p = 1/4$~: $\mathbb{E}[N] = 4$, et
$\mathbb{P}(N \le 4) = 1 - (3/4)^{4} = 1 - \dfrac{81}{256} = \dfrac{175}{256}
\approx 0{,}684.$

\section*{Partie V --- Partition équilibrée par programmation dynamique}

\Q Si $T$ est impaire, aucune part entière ne peut valoir $T/2$ (qui n'est pas entier)~; les deux
parts, de sommes entières égales, donneraient $T$ pair. Donc la réponse est \emph{non}. $\square$

\Q Un \emph{certificat} est la liste des indices d'un sous-ensemble candidat. La vérification
calcule sa somme et teste l'égalité à $C$, en temps $O(n)$ (polynomial). Le problème est donc
dans \textbf{NP}. (Il s'agit du problème \textsc{Subset-Sum}, NP-complet.)

\Q Un sous-ensemble des $i$ premiers éléments de somme $c$ ou bien n'utilise pas le $i$-ième
élément (donc c'est un sous-ensemble des $i-1$ premiers de somme $c$), ou bien l'utilise (donc
c'est un sous-ensemble des $i-1$ premiers de somme $c - S[i-1]$, licite si $c \ge S[i-1]$)~:
$$A[i][c] = A[i-1][c] \;\lor\; \bigl(c \ge S[i-1] \;\land\; A[i-1][c - S[i-1]]\bigr),$$
avec l'initialisation $A[0][0] = \text{vrai}$ et $A[0][c] = \text{faux}$ pour $c \ge 1$.

\Q
\begin{lstlisting}
def peut_partitionner(S):
    T = sum(S)
    if T % 2 == 1:
        return False
    C = T // 2
    dp = [False] * (C + 1)
    dp[0] = True
    for x in S:                       # traite un element a la fois
        for c in range(C, x - 1, -1): # parcours DECROISSANT
            if dp[c - x]:
                dp[c] = True
    return dp[C]
\end{lstlisting}
Le parcours \emph{décroissant} de $c$ garantit que \code{dp[c - x]} se réfère encore à l'état
\emph{avant} l'introduction de $x$~: on n'utilise donc jamais deux fois le même élément.
Complexité~: $O(nC)$ en temps et $O(C)$ en espace. Comme $C$ dépend de la \emph{valeur} des
données (et non seulement de leur nombre), la taille de $C$ est exponentielle en le nombre de bits
de codage~: l'algorithme est dit \emph{pseudo-polynomial}.

\Q \textbf{Invariant.} Après le traitement des $i$ premiers éléments, pour tout
$0 \le c \le C$, \code{dp[c]} est vrai si et seulement s'il existe un sous-ensemble des $i$
premiers éléments de somme $c$ (c'est-à-dire \code{dp[c]} $= A[i][c]$).

\noindent\emph{Initialisation} ($i=0$)~: seul \code{dp[0]} est vrai, ce qui correspond au
sous-ensemble vide~; conforme à $A[0]$. \emph{Hérédité}~: supposons l'invariant vrai après $i-1$
éléments et traitons $x = S[i-1]$. Pour chaque $c$ de $C$ à $x$, on met \code{dp[c]} à vrai
lorsque \code{dp[c - x]} est vrai. Le parcours décroissant assure que \code{dp[c - x]} vaut encore
$A[i-1][c-x]$~; combiné à la valeur préexistante $A[i-1][c]$, on obtient exactement
$A[i][c] = A[i-1][c] \lor (c \ge x \land A[i-1][c-x])$. L'invariant est donc vrai après $i$
éléments. \emph{Conclusion}~: après les $n$ éléments, \code{dp[C]} $= A[n][C]$, qui est vrai ssi
une partition équilibrée existe. \textbf{Terminaison}~: les deux boucles \code{for} sont finies.
$\square$

\medskip
\noindent\emph{Vérifications numériques}~: $\mathtt{peut\_partitionner}([1,1,5,7]) = \text{True}$
(part $[1,1,5]$ de somme $7$), $\mathtt{peut\_partitionner}([1,1,5,8]) = \text{False}$ (somme
$15$ impaire). Sur $2000$ jeux aléatoires, \code{peut_partitionner} coïncide avec une recherche
exhaustive.

\end{document}
